\documentclass[a4paper,12pt,tightenlines]{revtex4} 
\usepackage{epsfig}
\usepackage{graphicx}
\setlength{\textheight}{25.5cm}
\setlength{\textwidth}{16.5cm}
%\usepackage[none]{hyphenat} % THIS PACKAGE TO AVOID HYPHENATION, COMMENT IT IF YOU DON'T WANT TO
\usepackage{pslatex}


\begin{document}
\begin{center}
\begin{large}
{\bf Title of the Proposal}\\ 
\vspace{.2cm}
{\bf Experimental Setup}\\
\end{large}
\end{center}
\vspace{.4cm}
\begin{flushleft}
{\bf Spokesperson(s): Your Name}
\end{flushleft}
\thispagestyle{empty} 
\pagestyle{empty} 

%HERE YOU CAN START TO WRITE THE BODY OF THE PROPOSAL
\section{Scientific Motivation}

Template for the body of the experiment proposal. This is a proposal for an experiment 
to be performed with the TANDEM--ALPI--PIAVE ACCELERATOR of the National Laboratories 
of Legnaro. The length of the proposal~\cite{ref1} should not exceed 4 pages. 

You can use this template in a simple way by directly replacing this text with your text 
in any section (Titles, Sections, Captions, etc.). Figure~\ref{label_fig1} is an example 
of how to insert images. Put the figure caption below the corresponding figure.

\begin{figure}[h]
\begin{center}
%\includegraphics[width=\columnwidth]{fig1.eps}
\includegraphics[width=8cm]{fig1.jpg}
\end{center}
\caption{Insert here the caption of the figure. This is a long dummy text to check the editor.}
\label{label_fig1}
\end{figure}

\section{Proposed Experiment}

Insert here the description of the proposed experiment~\cite{ref2}. Specify which experimental setup will be used, special requirements, additional detectors. 

\begin{table}[h]
\caption{Insert your caption here. This is a long dummy text to check the editor.}
\begin{center}
\begin{tabular}{|c|c|c|c|}
\hline
\ E$_{c.m.}$ (MeV)  & ~$\sigma_{ER}$ (mb)  & E$_{c.m.}$ (MeV)  & ~$\sigma_{ER}$ (mb) \ \\ 
\hline
\ 85.36     &      0.0011$\pm$0.00079      & 95.97    &     116.2$\pm$1.8  \ \\
\ 85.84     &      0.0022$\pm$0.0010      & 96.93     &     132.2$\pm$2.9   \ \\
\hline
\end{tabular}
\label{table1}
\end{center}
\end{table}

You can also insert tables as shown in the above example (see table~\ref{table1}). Tables 
can be created directly (see table \ref{table1}) or inserted as an image. It is essential 
that the table content be easily 
readable: the character size must be adequate. The table caption has to be placed on top 
of the table. 

\section{Rate Estimates and Beam Time Request}

Estimates of the expected rates are required. The beam time request should specify: 
\begin{itemize}
\item beam: ion, intensity, energy, pulsed or continuous
\item target: isotope, thickness, backing
\item experimental setup:
\item beam time: xxx days
\end{itemize}

\begin{thebibliography}{999} 
\bibitem{ref1} A. First et al., Phys. Rev., C11 (2000) 1111.
\bibitem{ref2} B. Second et al., Nucl. Instr. and Meth., A111 (2000) 222.
\bibitem{ref3} C. Third et al., Eur. Phys. J., A222 (2000) 123.
\end{thebibliography}


\end{document}

